\documentclass[11pt]{article}
\usepackage[T1]{fontenc}
\usepackage{lmodern,amsmath,amssymb,amsthm,microtype,booktabs,xurl}
\usepackage[margin=26mm]{geometry}
\usepackage[hidelinks]{hyperref}
\newtheorem{theorem}{Theorem}
\newtheorem{proposition}[theorem]{Proposition}
\newtheorem{lemma}[theorem]{Lemma}
\newtheorem{corollary}[theorem]{Corollary}
\newcommand{\Z}{\mathbb Z}\newcommand{\Q}{\mathbb Q}
\newcommand{\workbenchurl}{https://github.com/KokunoYumeto/zeta-and-l-functions-workbench}
\newcommand{\recordurl}{https://github.com/KokunoYumeto/mathematics-workbenches}
\title{Strict seven-adic Selmer groups and universal norms in a cubic tower, after Fukuda and Rubin}
\author{}\date{24 November 2026}
\begin{document}\maketitle
\begin{abstract}
Fukuda's stability theorem and Rubin's class-field exact sequence connect finite class groups with cyclotomic Iwasawa modules. We certify that the seven-primary class groups of the real cubic field generated by the cube root of 269 and of its first cyclotomic seven-layer both have order seven. Every subsequent class group has the same seven-part. Over the cubic splitting field this determines both faithful character components as finite Iwasawa modules and proves that every matched principal cyclotomic weight has strict Selmer length exactly one. The one-prime ramified module and the corresponding local-unit quotient, in contrast, have a linear characteristic polynomial with a constant coefficient of valuation one. We give a full proof from capitulation, universal norms, and two exact local unit matrices. The values of this polynomial are unbounded along positive integer weights, so a bare characteristic-generator formula for strict lengths is false. The seven-adic contribution to an all-prime arithmetic sum is explicit; the remaining bound needed for the cubic Riemann hypothesis is still open. The mathematics and text were produced by GPT-6 Astra in ChatGPT Pro mode; no human has checked this text.
\end{abstract}
\section{Exact stability at seven and all cyclotomic characters}
\label{seven:section}

Fukuda's stability theorem reduces a class-group calculation in an entire
$\mathbb Z_p$-extension to two successive finite layers
\cite[Theorem~1(1), p.~264]{sevenFukuda1994}.
We give its short finite-level proof in the one-prime situation needed here
and certify the two arithmetic inputs. This gives an exact finite Iwasawa
module, including its finite submodule, and hence determines every strict
cyclotomic specialization.

Put
\[
 F=\mathbb Q(\alpha),\quad \alpha^3=269,\qquad
 K=\mathbb Q(\omega),\quad E=KF,
\]
where $\alpha$ is the real root and $\omega$ is a primitive cube root of unity.
Let $Q_n$ be the degree-$7^n$ layer of the cyclotomic
$\mathbb Z_7$-extension of $\mathbb Q$, and set
$F_n=FQ_n$, $K_n=KQ_n$, and $E_n=EQ_n$.
Write $A(k)=\operatorname{Cl}(k)\{7\}$, and
\[
 \Delta=\operatorname{Gal}(E/K)=\langle\delta\rangle,\qquad
 \Gamma=\operatorname{Gal}(E_\infty/E).
\]
Both faithful characters of $\Delta$ take values in $\mu_3\subset\mathbb Z_7$.
Choose $\gamma\in\Gamma$ with
$\langle\varepsilon_7\rangle(\gamma)=8$, and put
$T=\gamma-1$ and $\Lambda=\mathbb Z_7[[T]]$.

For a finite extension $V/\mathbb Q_7$, with integer ring $\mathcal O$, and
a continuous character $\kappa:\Gamma\longrightarrow\mathcal O^\times$,
let
\[
 D_{\chi,\kappa}=(V/\mathcal O)(\chi\kappa).
\]
The characters here are inflated through
$\operatorname{Gal}(E_\infty/K)=\Delta\times\Gamma$.
The strict Selmer group means
\[
 \operatorname{Sel}_{\rm str}(K,D)=
 \ker\left(
 H^1(K,D)\longrightarrow
 \prod_{v\mid7}H^1(K_v,D)
 \times\prod_{v\nmid7}H^1(I_v,D)\right).
 \tag{S1}\label{seven:strict-definition}
\]
The second product is over finite places; imposing the usual condition at
infinite places has no effect on seven-primary cohomology.

\begin{theorem}\label{seven:main}
For every $n\ge0$,
\[
 A(F_n)\simeq\mathbb Z/7\mathbb Z.
 \tag{S2}\label{seven:class-stability}
\]
All norms between these groups are isomorphisms and all cyclotomic actions
are trivial. For either faithful $\chi$ of $\Delta$,
\[
 \left(\varprojlim_n A(E_n)\right)^\chi
 \simeq\Lambda/(7,T).
 \tag{S3}\label{seven:iwasawa}
\]
For every continuous $\kappa$ as above, with $u=\kappa(\gamma)$,
\[
 \operatorname{Sel}_{\rm str}(K,D_{\chi,\kappa})
 \simeq\mathcal O/(7,u-1).
 \tag{S4}\label{seven:all-characters}
\]
These assertions are unconditional.
\end{theorem}

The finite certificate uses exact algebraic integers, a matrix of
residue characters, and bounded convergent series. A list of units alone
does not certify completeness of a class group or a unit basis. The certificate proves precisely the seven-part required in
\eqref{seven:class-stability}; it does not need a determination of every
prime divisor of either class number.

\subsection{The one-prime module and its weight values}
\label{seven:universal-norms}

The finite strict module of Theorem~\ref{seven:main} does not coincide
with the Iwasawa module in which ramification is allowed at one of the
two primes above seven. We determine the latter module and its unit
quotient. The distinction comes from capitulation: the faithful class
groups remain of order seven under norms, but their extension maps
vanish, and the resulting defect in norms of global units persists
throughout the tower.

Fix a prime $\mathfrak p\mid7$ of $K=\mathbb Q(\omega)$ and a faithful
character $\chi$ of $\Delta=\operatorname{Gal}(E/K)$. Choose
$\gamma\in\Gamma=\operatorname{Gal}(E_\infty/E)$ with
$\langle\varepsilon_7\rangle(\gamma)=8$, and put
\[
 \Lambda=\mathbb Z_7[[T]],\qquad T=\gamma-1,\qquad
 \omega_n(T)=(1+T)^{7^n}-1.
\]
Let $L_n$ be the completion of $E_n$ at its unique prime over
$\mathfrak p$, and define
\[
 U_n=U^1(L_n),\qquad
 \mathcal E_n=\mathbb Z_7\otimes_{\mathbb Z}\mathcal O_{E_n}^{\times},
 \qquad
 U_\infty=\varprojlim_N U_n,\qquad
 \mathcal E_\infty=\varprojlim_N\mathcal E_n.
\]
The localization map $\ell_n:\mathcal E_n\to U_n$ uses the principal-unit
projection. Its image $\overline{\mathcal E}_n$ is the closure of the
global units which are principal at this prime. Write
$\overline{\mathcal E}_\infty=\varprojlim_N\overline{\mathcal E}_n$ and
\[
 V_{\mathfrak p,\chi}
   =(U_\infty/\overline{\mathcal E}_\infty)^\chi.
\]
Finally, let $X_{\mathfrak p,\infty}$ be the Galois group of the maximal
abelian pro-seven extension of $E_\infty$ unramified outside the primes
over $\mathfrak p$. These are the one-prime local-unit and Galois modules
of \cite[printed p.~36]{Rubin}.

\begin{theorem}\label{seven:one-prime-linear}
For each choice of $\mathfrak p$ and faithful $\chi$, there is
$a_{\mathfrak p,\chi}\in7\mathbb Z_7^\times$ such that
\[
 U_\infty^\chi\simeq\Lambda,\qquad
 \mathcal E_\infty^\chi\simeq\Lambda,\qquad
 V_{\mathfrak p,\chi}\simeq
 X_{\mathfrak p,\infty}^\chi
 \simeq\Lambda/(T-a_{\mathfrak p,\chi}).
\]
The class-field exact sequence is
\begin{equation*}\tag{U1}\label{seven:one-prime-exact}
 0\longrightarrow V_{\mathfrak p,\chi}
 \longrightarrow X_{\mathfrak p,\infty}^\chi
 \longrightarrow M^\chi
 \longrightarrow0,
\end{equation*}
where $M^\chi=\Lambda/(7,T)$, and its first map identifies
$V_{\mathfrak p,\chi}$ with $7X_{\mathfrak p,\infty}^\chi$.
In particular the one-prime characteristic ideal has
$\mu=0$, $\lambda=1$, and is not the unit ideal.

Put
\[
 e_{\mathfrak p,\chi}^{*}
 =-\frac{\log(1+a_{\mathfrak p,\chi})}{\log8}
 \in\mathbb Z_7^\times.
\]
For every generator $f_{\mathfrak p,\chi}$ of
$\operatorname{char}_{\Lambda}X_{\mathfrak p,\infty}^\chi$ and every
$e\in\mathbb Z_7$,
\begin{equation*}\tag{U2}\label{seven:one-prime-weight-valuation}
 v_7\!\left(f_{\mathfrak p,\chi}(8^{-e}-1)\right)
 =1+v_7(e-e_{\mathfrak p,\chi}^{*}),
\end{equation*}
with both sides interpreted as infinity at
$e=e_{\mathfrak p,\chi}^{*}$. These valuations are unbounded on the
positive integer weights, whereas their strict Selmer lengths in
Theorem~\ref{seven:main} are all one.
\end{theorem}


\subsection{The exact seven-adic contribution to the arithmetic sum}

Let $T\ge2$ be an integer and $N=T-1$. Let $\xi_7$ be either faithful character with values in $\mu_6$ obtained from $\operatorname{Gal}(\Q(\zeta_7)/\Q)$. For each $e\in\{1,\ldots,N\}$, let $a(e)\in\{0,\ldots,5\}$ be the unique exponent satisfying
\[
 \overline\xi_7^{\,a(e)}\overline\varepsilon_7^{-e}=1.
\]
The matched family consists of the two coefficients
\[
 (V/\mathcal O)(\chi^j\xi_7^{a(e)}\varepsilon_7^{-e}),\qquad j=1,2,
\]
at each weight. The strict conditions are those in the theorem: complete local triviality at both primes over seven and unramifiedness at every other finite prime. Put $b=\operatorname{rank}_{\Z_7}\mathcal O$.

\begin{corollary}\label{cor:seven-budget}
Let $J$ be any finite family of nonnegative normalized arithmetic terms containing these $2N$ matched terms, each counted once. Write $\mathcal R_J$ for their sum and $\mathcal R_{\rm rest}$ for the sum after removing the matched terms. Then
\[
 \mathcal R_J=\mathcal R_{\rm rest}+2N\log7.
 \tag{B1}
\]
If $\mathcal R^{[m]}$ instead leaves the matched pair at weight $m$, then
\[
 \mathcal R^{[m]}=\mathcal R_{\rm rest}+2\log7,
 \qquad
 \max_{1\le m\le N}\mathcal R^{[m]}
       =\mathcal R_{\rm rest}+2\log7.
 \tag{B2}
\]
Suppose a comparison on the same datum satisfies
\[
 |I-\mathcal R_J|\le E+C\mathcal P,
 \qquad \mathcal P\ge T\log2.
\]
Then
\[
 |I-\mathcal R_{\rm rest}|
 \le E+C\mathcal P+2(T-1)\log7
 < E+(C+6)\mathcal P.
 \tag{B3}
\]
For every real threshold $B$ the exact equivalence is
\[
 \mathcal R_J\le B
 \quad\Longleftrightarrow\quad
 \mathcal R_{\rm rest}\le B-2(T-1)\log7.
 \tag{B4}
\]
\end{corollary}

The abstract residue-class profiles below are auxiliary comparisons:
they do not assert exceptional strict weights in this tower.
Theorem~\ref{th:profile} and Propositions~\ref{prop:peaks}--\ref{prop:sum}
describe those conditional profiles, while
Lemma~\ref{lem:euler-specialization} supplies the general algebraic
specialization correction.

\section*{Proof overview}

We first compute the quotient of each class number by the index of an
explicit group of units. A universal regulator lower bound makes the
possible denominator finite. Rigorous intervals isolate the quotient as
the integer seven, while independent residue characters prove that the
unit index is prime to seven. In the cyclotomic tower of $F$, exactly one
prime ramifies. Class field theory therefore identifies class-group
coinvariants with the lower class group. Equality at the first two levels
and Nakayama's lemma force stability at every level. The two faithful
cubic-character parts over $E$ are recovered from the conjugation-fixed
part. Finally, inflation and restriction preserve both strict local
conditions because the faithful cubic character has no invariants.
The resulting specialization of the actual finite module is
\eqref{seven:all-characters}. The unit calculation then uses the vanishing of extension maps on these class groups. Hilbert 90 and the faithful Herbrand quotient give a global unit norm index of seven at every pair of levels. Passing to inverse limits preserves this defect. A direct calculation in the first local layer gives index seven for the localization of global units; together these indices force the local-unit quotient to have a linear characteristic polynomial. Rubin's theorem excluding finite submodules identifies the one-prime Galois module and makes the weight valuation formula exact.

\subsection{A finite certification principle}
\label{seven:finite-principle}

\begin{lemma}\label{seven:isolation}
Let $k$ be a number field with regulator $R_k>1/4$. Let $U_0$ contain all
roots of unity of $k$ and be generated modulo torsion by $r$ independent
units, where $r$ is the unit rank. Put
$m=[\mathcal O_k^\times:U_0]$ and $R_0=mR_k$.
Suppose $B$ is an integer with $m<B$ and an interval calculation proves
\[
 \left|\frac{h_k}{m}-7\right|<\frac1B.
 \tag{S5}\label{seven:isolation-bound}
\]
Then $h_k=7m$. If $r$ homomorphisms
$\mathcal O_k^\times\longrightarrow\mathbb F_7$ have a nonsingular
evaluation matrix on the specified units, then $7\nmid m$, and
$A(k)\simeq\mathbb Z/7\mathbb Z$.
\end{lemma}

\begin{proof}
The relation $R_0=mR_k$ follows by comparing the covolumes of the logarithmic
unit lattices. Since $h_k-7m$ is an integer of absolute value less than
$m/B<1$, it vanishes. If $7$ divided $m$, the lattice inclusion
$U_0/\mu(k)\subset\mathcal O_k^\times/\mu(k)$ would have a singular
reduction modulo seven. Equivalently, there would be a product of the
specified units, with exponents not all divisible by seven, which is a
seventh power of a unit. Every displayed homomorphism would kill this
product, contradicting the nonsingular matrix. Thus
$v_7(h_k)=1$, as asserted.
\end{proof}

The regulator bound used here is Friedman's
\cite[Theorem~B, pp.~599--600]{sevenFriedman1989}: $R_k>1/4$ outside three
exceptional totally complex sextic fields. Both fields used below have real
embeddings and degrees three and twenty-one, so this exception is irrelevant.
In particular, one may take $B=\lceil4R_0\rceil$, with an upper enclosure of
$R_0$ used when calculating the ceiling.

For a rational prime $\ell\equiv1\pmod7$ and a prime
$\mathfrak l$ of $k$ with residue field $\mathbb F_\ell$, choose
$z\in\mathbb F_\ell^\times$ of order seven. The character
\[
 u^{(\ell-1)/7}\equiv z^{\,\ell_{\mathfrak l}(u)}
 \pmod{\mathfrak l},\qquad
 \ell_{\mathfrak l}(u)\in\mathbb F_7
 \tag{S6}\label{seven:residue-character}
\]
is a homomorphism on all units. This supplies the homomorphisms in the lemma.
The fields here have a real embedding and odd degree, so their roots of
unity are exactly $\{\pm1\}$.

\subsection{The zeta residues as bounded finite sums}
\label{seven:analytic}

Set $N=217083=3\cdot269^2$, and let $\rho$ be the standard two-dimensional
representation of $\operatorname{Gal}(E/\mathbb Q)\simeq S_3$. Then
\[
 L_0(s)=L(s,\rho)=\frac{\zeta_F(s)}{\zeta(s)},\qquad
 D_F=-N.
 \tag{S7}\label{seven:base-L}
\]
For completeness, the discriminant exponent at $269$ is two by tame total
ramification. At three, $v_3((-1)^3-269)=3>2v_3(3)=2$, so Hensel's
lemma gives a root in $\mathbb Q_3$ and
$F\otimes\mathbb Q_3=\mathbb Q_3\times\mathbb Q_3(\omega)$.
The exponent at three is therefore one. These observations also identify
the ramified Euler factors below.

Let $\psi$ be the order-seven character modulo $49$ with
$\psi(3)=\zeta_7$, and put $L_j=L(s,\rho\otimes\psi^j)$.
For $1\le j\le6$, the character $\psi^j$ is even and primitive of conductor
$49$: $3$ has order $42$ modulo $49$, and $-1=3^{21}$.
Its conductor cannot divide seven because
$|(\mathbb Z/7\mathbb Z)^\times|=6$.
Since $\rho$ is unramified at seven, twisting its two unramified dimensions
gives conductor exponent four there. Thus
\[
 N_0=N,\quad N_j=49^2N\ (j\ne0),\qquad
 \Lambda_j(s)=N_j^{s/2}\Gamma_{\mathbb C}(s)L_j(s),
 \quad
 \Gamma_{\mathbb C}(s)=2(2\pi)^{-s}\Gamma(s).
 \tag{S8}\label{seven:conductors}
\]
The Dedekind functional equation gives
$\Lambda_0(s)=\Lambda_0(1-s)$
\cite[Remark~18.14]{sevenSutherland2015}.
The representation is induced from a nontrivial cubic Hecke character of
$K$ of conductor $(269)$, whose conjugate is its distinct inverse.
Hecke's theta construction therefore gives the cuspidal weight-one form
$f(z)=\sum_{n\ge1}a_ne^{2\pi inz}$ of level $N$ and nebentypus
$\eta=\left(\frac{-3}{\,\cdot\,}\right)$; see
\cite[Theorem~5]{sevenSimard2015} for the theta transformation and
\cite[\S3]{sevenDarmonLauderRotger2015} for the weight-one induced
representation and cuspidality criterion.

Here is a derivation of the twisting sign, to fix its orientation.
For a primitive even character $\xi\bmod M$, $M=49$, set
$\tau(\xi)=\sum_{a\bmod M}\xi(a)e^{2\pi ia/M}$ and
\[
 A_a=\begin{pmatrix}1&a/M\\0&1\end{pmatrix},\qquad
 W_Q=\begin{pmatrix}0&-1\\Q&0\end{pmatrix}.
\]
With the normalized weight-one slash operator, finite Fourier inversion gives
\[ f_\xi=\tau(\bar\xi)^{-1}\sum_a\bar\xi(a)f|A_a. \]
Choose $b$ with $Nab\equiv-1\pmod M$. Direct multiplication gives
\[
 A_aW_{NM^2}=
 \begin{pmatrix}(1+Nab)/M&a\\Nb&M\end{pmatrix}M W_NA_b.
 \tag{S9}\label{seven:fricke-matrix}
\]
The first matrix has determinant one, lies in $\Gamma_0(N)$, and has
lower-right entry $M$, with $\eta(M)=1$.
Also $\bar\xi(a)=\xi(N)\xi(b)$.
The base functional equation gives $f|W_N=-if$ by Mellin inversion.
Using $\tau(\xi)\tau(\bar\xi)=M$ consequently gives
\[
 \Lambda_j(s)=\varepsilon_j\Lambda_{7-j}(1-s),\qquad
 \varepsilon_j=\psi^j(N)\frac{\tau(\psi^j)^2}{49}.
 \tag{S10}\label{seven:epsilon}
\]
There is a short exact verification of the Gauss sums.
The relation $3^6=1-7\bmod49$ gives
$\psi(1+7b)=\zeta_7^b$. Grouping terms as $a+7t$ shows that only
$a\equiv-j\pmod7$ survives in $\tau(\psi^j)$.
Since $N\equiv3^{33}\pmod{49}$, the results are
\[
\begin{array}{c|rrrrrr}
 j&1&2&3&4&5&6\\ \hline
 \tau(\psi^j)/7=\zeta_{49}^{\,b_j}
   &48&19&18&31&30&1\\
 \varepsilon_j=\zeta_{49}^{\,c_j}
   &33&10&43&6&39&16
\end{array}
 \tag{S11}\label{seven:gauss-table}
\]
where the entries in the last two rows are the exponents $b_j,c_j$.

The coefficients $a_n$ are multiplicative and have the following complete
prime-power specification, including $a_{p^0}=1$:
\[
\begin{array}{c|l}
 p& a_{p^k}\quad(k\ge1)\\ \hline
269&0\\
3&1\\
p\equiv2\pmod3,\ p\ne269
   &1\text{ for even }k,\quad0\text{ for odd }k\\
p\equiv1\pmod3,\ 269^{(p-1)/3}\equiv1\pmod p&k+1\\
p\equiv1\pmod3,\ 269^{(p-1)/3}\not\equiv1\pmod p
   &1,-1,0\text{ for }k\equiv0,1,2\pmod3 .
\end{array}
 \tag{S12}\label{seven:coefficients}
\]
These are the Euler factors for the identity, transposition, and
three-cycle classes, together with the inertia invariants at the two
ramified primes. In particular, $|a_n|\le d(n)\le n$.

Put $\lambda_j=2\pi/\sqrt{N_j}$.
Splitting the Mellin integral at one and using
\eqref{seven:epsilon} yields, with
$E_1(x)=\int_x^\infty e^{-t}\,dt/t$,
\[
 L_j(1)=
 \sum_{n\ge1}\frac{a_n\psi^j(n)e^{-\lambda_jn}}{n}
 +\varepsilon_j \lambda_j
 \sum_{n\ge1}a_n\overline{\psi^j(n)}E_1(\lambda_jn).
 \tag{S13}\label{seven:L1-formula}
\]
For $j=0$ the character in this formula is identically one, including at
multiples of seven, and $\varepsilon_0=1$.
Indeed, the theta transformation is
$\Theta_j(t)=\varepsilon_jt^{-1}\Theta_{7-j}(1/t)$; after changing
$t$ to $1/t$ in the lower Mellin integral, its value at $s=1$ is precisely
the second term in \eqref{seven:L1-formula}.
Because $0<E_1(x)\le e^{-x}/x$, deleting all terms $n>X$ in both sums
has joint absolute error at most
\[
 {\cal T}(c,X)=\frac{2e^{-c(X+1)}}{1-e^{-c}}.
 \tag{S14}\label{seven:tail}
\]
This estimate uses only $|a_n|\le n$ and a geometric series.

The degree-seven field $Q_1$ has
$\zeta_{Q_1}(s)=\prod_{j=0}^6L(s,\psi^j)$.
Here $\psi^0$ means the trivial character of conductor one, whose factor is
$\zeta(s)$, so its Euler factor at seven is retained.
For $j\ne0$ its factors at one are given by the finite identity
\[
 L(1,\psi^j)=
 -\frac{\tau(\psi^j)}{49}
 \sum_{\substack{1\le a<49\\7\nmid a}}
       \overline{\psi^j(a)}
       \log\left(2\sin\frac{\pi a}{49}\right).
 \tag{S15}\label{seven:dirichlet-residue}
\]
This follows by Fourier expansion of $-\log(2\sin\pi x)$ and the
Gauss-sum identity; evenness combines the two Fourier directions.
Pairing conjugate characters makes the resulting residue real and positive.
Finally, the permutation representation for $F_1/\mathbb Q$ is
$(1+\rho)\otimes\bigoplus_{j=0}^6\psi^j$, so Artin formalism gives
\[
 \operatorname*{res}_{s=1}\zeta_{F_1}(s)
 =\operatorname*{res}_{s=1}\zeta_{Q_1}(s)
       \prod_{j=0}^6L_j(1).
 \tag{S16}\label{seven:residue-factorization}
\]
This identity includes every Euler factor, including those at seven.

\subsection{Explicit units and rigorous interval isolation}
\label{seven:certification}

Let
\[
 q(Y)=Y^7-21Y^5-21Y^4+91Y^3+112Y^2-84Y-97,
 \qquad
 \beta=\sum_{a\in\{1,18,19,30,31,48\}}\zeta_{49}^{a}.
 \tag{S17}\label{seven:period}
\]
The set of exponents is the subgroup of order six in
$(\mathbb Z/49\mathbb Z)^\times$. Thus $\mathbb Q(\beta)=Q_1$.
The polynomial $q$ is irreducible and $q(\beta)=0$ by direct arithmetic
modulo $\Phi_{49}$.
Write $\theta=\alpha-\beta$. The polynomial $p$ of degree twenty-one and
the rational function $P/Q$ printed in the finite data below satisfy
\[
 p(\theta)=0,\qquad \alpha=P(\theta)/Q(\theta),\qquad
 \alpha^3=269,\qquad q(\alpha-\theta)=0.
 \tag{S18}\label{seven:field-identity}
\]
Irreducibility of $p$ and $q$, or the relatively prime degrees three and
seven together with these identities, identifies this field as $F_1$.
The automorphism $\tau$ fixes $\alpha$ and is
$\zeta_{49}\mapsto\zeta_{49}^3$ on $Q_1$; its action on $\beta$ is the
degree-six rational polynomial printed below.

The first six units $u_1,\ldots,u_6$ lie in $\mathbb Q(\beta)$ and are
given by their seven coefficients and denominators.
The seventh is
\[
 u_7=-\varepsilon,\qquad
 \varepsilon=270-739\alpha+108\alpha^2.
 \tag{S19}\label{seven:base-unit}
\]
For the explicitly printed $\eta$, the remaining units are
$u_{8+i}=\tau^i\eta$ for $0\le i\le5$.
Exact rational polynomial arithmetic verifies that all thirteen units are
integral and have absolute norm one. One can check integrality either in
the certified maximal order or from the monic integral characteristic
polynomial of multiplication by the element. The discriminant and
signature are
\[
 D_{F_1}=-N^7\,7^{36},\qquad (r_1,r_2)=(7,7).
 \tag{S20}\label{seven:disc-signature}
\]
Independently, $D_{Q_1}=49^6=7^{12}$ by the conductor--discriminant
formula. Its discriminant is coprime to $D_F$, so the product of the
integral bases is a maximal order in the compositum and gives
\eqref{seven:disc-signature}.

Here are complete numerical specifications for the interval calculation.
Generate the integer coefficients in \eqref{seven:coefficients} by a
prime-factor sieve. In \eqref{seven:L1-formula} take $X=14000$ for $j=0$
and $X=550000$ for $j\ne0$, and evaluate all terms by real and complex
ball arithmetic at precision $300$ bits.
The bounds \eqref{seven:tail} are, respectively,
$1.496\cdot10^{-80}$ and $1.328\cdot10^{-62}$.
The resulting enclosures, rounded outward, include
\[
\begin{split}
 L_0(1)&\in
 [\,1.6976757232625543059851766859761711810674
       \ \pm 10^{-40}\,],\\
 \prod_{j=0}^6L_j(1)&\in
 [\,0.6494049966247962459428741324547233673973
       \ \pm 10^{-40}\,].
\end{split}
 \tag{S21}\label{seven:analytic-intervals}
\]
For the regulator, isolate the seven disjoint real roots of $q$ with
complex ball arithmetic, take
$\alpha_0=\exp(\log269/3)$ and
$\alpha_1=\alpha_0(-1+i\sqrt3)/2$, and evaluate the thirteen units at
$(\alpha_0,\beta_j)$ and $(\alpha_1,\beta_j)$.
Form the logarithmic matrix with entries $\log|\sigma(u_i)|$ for real
embeddings and $2\log|\sigma(u_i)|$ for complex pairs, delete one real
row, and take the absolute determinant.
At $2048$ bits this gives the regulator $R_0$ of the displayed unit group
with over $600$ bits of relative accuracy. Arb's enclosures, including
the root isolation and determinant operations, are used throughout
\cite{sevenJohansson2017}.
Rounded outward at lower precision, the relevant values are
\[
\begin{split}
 R_0&\in
 [\,8.7214139234209020456623583408322697739521
       \ \pm10^{-40}\,]\,10^{24},\\
 \operatorname*{res}_{s=1}\zeta_{Q_1}(s)
   &\in[\,0.2996769180948235445839119514335675929714
       \ \pm10^{-40}\,].
\end{split}
 \tag{S22}\label{seven:regulator-intervals}
\]
The positive determinant proves independence of the units.

The analytic class-number formula, in the convention just used, is
\[
 \operatorname*{res}_{s=1}\zeta_k(s)
  =\frac{2^{r_1}(2\pi)^{r_2}h_kR_k}
         {|\mu(k)|\sqrt{|D_k|}}
 \tag{S23}\label{seven:class-formula}
\]
\cite[Theorem~18.13]{sevenSutherland2015}.
If $m_1$ is the index of the displayed thirteen units and $\{\pm1\}$,
\eqref{seven:residue-factorization} and the unrounded ball evaluations give
\[
 \frac{h_{F_1}}{m_1}
 =\frac{\sqrt{N^7\,7^{36}}}
        {2^{13}\pi^7R_0}
   \operatorname*{res}_{s=1}\zeta_{Q_1}(s)\prod_{j=0}^6L_j(1)
 \in[\,7\pm1.11\cdot10^{-59}\,].
 \tag{S24}\label{seven:quotient-interval}
\]
The stated precision and cutoffs, rather than substitution of the rounded
displayed decimals, specify this computation.
Friedman's bound gives
\[
 m_1<B=34885655693683608182649434,
 \qquad
 B\left|h_{F_1}/m_1-7\right|<3.88\cdot10^{-34}<1.
 \tag{S25}\label{seven:index-bound}
\]
The thirteen residue primes and their full evaluation matrix are printed
below. Their determinant is $1$ in $\mathbb F_7$.
Lemma~\ref{seven:isolation} proves $v_7(h_{F_1})=1$.

The same short certificate supplies the base layer without assuming a
class-number search. For \eqref{seven:base-unit}, exact arithmetic gives
$N_{F/\mathbb Q}(\varepsilon)=1$ and
\[
 0<R_\varepsilon=|\log|\varepsilon||<18,\qquad
 \frac{\sqrt N L_0(1)}{2\pi R_\varepsilon}
 \in[\,7\pm4.0\cdot10^{-68}\,].
 \tag{S26}\label{seven:base-quotient}
\]
Thus $m_0=[\mathcal O_F^\times:\langle-1,\varepsilon\rangle]<72$ and
$h_F=7m_0$.
At the prime $491$ the map $\alpha\mapsto444$ gives
$\varepsilon\mapsto88$, while
\[
 138^7\equiv1,\quad138\not\equiv1,\qquad
 88^{70}\equiv223\equiv138^4\not\equiv1\pmod{491}.
 \tag{S27}\label{seven:base-saturation}
\]
This nonzero residue character proves $7\nmid m_0$.
Consequently
\[
 A(F)=A(F_1)=\mathbb Z/7\mathbb Z.
 \tag{S28}\label{seven:finite-inputs}
\]
All integers, unit expressions, coefficient rules, cutoffs and tail bounds
needed for these two assertions have now been specified.

\subsection{Norm control and stabilization}
\label{seven:stability-proof}

The cube classes in $\mathbb F_7$ are $0,1,6$, whereas $269\equiv3$.
Thus $X^3-269$ is irreducible modulo seven, and seven is unramified and inert
in $F$. The local extension at this unique prime is the unramified cubic
extension of $\mathbb Q_7$.
Each cyclotomic extension $Q_n/Q_m$ is totally ramified at seven and
unramified everywhere else. Base change by the unramified cubic local
field preserves this total ramification and its degree. Therefore every
$F_n/F_m$ has exactly one ramified finite prime, which is totally ramified.

\begin{lemma}\label{seven:norm-control}
Let $L/k$ be cyclic of degree $7^r$, with exactly one ramified finite prime,
which is totally ramified. Then the ideal norm induces an isomorphism
\[
 A(L)_{\operatorname{Gal}(L/k)}\ \xrightarrow{\ \sim\ }\ A(k).
 \tag{S29}\label{seven:coinvariant-control}
\]
\end{lemma}

\begin{proof}
Let $H$ be the Hilbert seven-class field of $L$. It is Galois over $k$,
since its defining maximality property is invariant under every
automorphism of $L/k$. Set
$G=\operatorname{Gal}(H/k)$ and
$A=\operatorname{Gal}(H/L)=A(L)$.
An inertia group $I$ above the sole ramified prime has order $7^r$:
$H/L$ is unramified and $L/k$ is totally ramified.
It therefore maps isomorphically to $\operatorname{Gal}(L/k)$, giving
$G=A\rtimes I$.
If $\sigma$ generates $I$, then $[G,G]=(\sigma-1)A$.
Killing $I$ and $[G,G]$ kills all inertia groups and gives the maximal
abelian unramified quotient of $G$, namely $A/(\sigma-1)A$.
The Hilbert seven-class field $H_k$ of $k$ is contained in $H$, since
$H_kL/L$ is unramified, abelian and of seven-power degree. Conversely,
every abelian unramified subextension of $H/k$ lies in $H_k$.
The quotient is therefore $A(k)$.
Under reciprocity this restriction map is the ideal norm, proving the
claim. There is no archimedean inertia of odd order to consider.
\end{proof}

This is the one-prime case of the usual class-field-theoretic norm
calculation; compare
\cite[Corollaries~1.3.9 and 1.3.12]{sevenSharifi}.
The proof includes the finite control map, so that no unidentified
finite kernel is carried into the next argument.

Fix $n\ge1$, put $A_n=A(F_n)$, and regard it as a module over
$R_n=\mathbb Z_7[C_{7^n}]$.
This ring is local with maximal ideal $(7,T)$, because
$R_n/7R_n=\mathbb F_7[T]/(T^{7^n})$.
Lemma~\ref{seven:norm-control} gives
\[
 A_n/TA_n\simeq A(F),\qquad
 A_n/\omega_1(T)A_n\simeq A(F_1),\qquad
 \omega_1(T)=(1+T)^7-1=T\nu(T),
 \tag{S30}\label{seven:two-quotients}
\]
where
$\nu(T)=7+21T+35T^2+35T^3+21T^4+7T^5+T^6$.
For $n=1$ the second quotient is simply $A_1$, so the statement remains
valid. Both quotients have order seven by \eqref{seven:finite-inputs}.
Since $\omega_1A_n\subseteq TA_n$, equality of their indices implies
\[
 TA_n=\omega_1A_n=\nu(T)\,TA_n.
\]
But $\nu(T)\in(7,T)$, so Nakayama's lemma applied to the finite
$R_n$-module $TA_n$ gives $TA_n=0$.
The first quotient in \eqref{seven:two-quotients} now proves
\eqref{seven:class-stability}.
Applying Lemma~\ref{seven:norm-control} to every $F_n/F_m$ shows that all
norms are surjective, hence isomorphisms between groups of order seven.
The resulting inverse limit is $\Lambda/(7,T)$.
This also proves directly the needed case of Fukuda's theorem
\cite[Theorem~1(1)]{sevenFukuda1994}.

The same norm lemma applied to $Q_n/\mathbb Q$, followed by Nakayama,
gives
\[
 A(Q_n)=0 \qquad(n\ge0).
 \tag{S31}\label{seven:rational-layer}
\]
Indeed its coinvariants are $A(\mathbb Q)=0$.

\subsection{The two faithful cubic-character components}
\label{seven:character-components}

Let $c$ be complex conjugation. Then
$\operatorname{Gal}(E_n/Q_n)=S_3$ with
$c\delta c=\delta^{-1}$, and $E_n^{\langle c\rangle}=F_n$.
Set $C_n=A(E_n)$.
For an extension of degree $d$ prime to seven, extension and norm of ideals
satisfy
$Nj=d$ and $jN=\sum_g g$.
Since $d$ acts invertibly on seven-primary groups, extension identifies the
lower class group with the invariant subgroup, without a ramification
restriction. Thus
\[
 C_n^{c=1}\simeq A(F_n),\qquad
 C_n^\Delta\simeq A(K_n),\qquad
 (C_n^\Delta)^{c=1}\simeq A(Q_n)=0.
 \tag{S32}\label{seven:prime-to-seven-descent}
\]
For $\varphi\in\{1,\chi,\chi^{-1}\}$ put
$e_\varphi=\frac13\sum_{d\in\Delta}\varphi(d)^{-1}d$.
The idempotents split
$C_n=C_n^1\oplus C_n^\chi\oplus C_n^{\chi^{-1}}$, and $c$ exchanges the
last two summands.
Since the $c$-fixed part of $C_n^1$ vanishes, projection gives an
isomorphism
\[
 e_\chi:C_n^{c=1}\ \xrightarrow{\ \sim\ }\ C_n^\chi,
 \qquad x\longmapsto e_\chi x,
 \qquad y\longmapsto y+cy
 \tag{S33}\label{seven:projection}
\]
with the displayed inverse.
Extension of ideals, the idempotents, and these projections commute with
the cyclotomic action and with the norm maps in the compositum tower.
Equations \eqref{seven:class-stability} and \eqref{seven:projection}
therefore prove \eqref{seven:iwasawa} for both faithful characters.

Only $(C_n^1)^{c=1}=0$ was used; deleting $C_n^1$ itself would be false.
For example, $K$ has class number one by its Minkowski bound
$(2/\pi)\sqrt3<2$. In $K_n/K$ the two primes above seven are totally
ramified, and each unit of $K$ is the norm of a unit, since the norm on
$\mu_6$ is raising to $7^n$.
The cyclic ambiguous class-number formula consequently gives
$|A(K_n)^{\operatorname{Gal}(K_n/K)}|=7^n$.
This growth lies in the conjugation-odd part and does not affect
\eqref{seven:projection}; the relevant cyclic class-field exact sequence
is given in \cite[Theorem~1.3.16]{sevenSharifi}.

\subsection{Strict local control and specialization}
\label{seven:control}

Seven splits into two primes in $K$, and $E/K$ is unramified and inert of
degree three at each of them. Hence each prime over seven in $E_n$ is fixed
by $\Delta$.
If
\[
 C_n'=\operatorname{Cl}(\mathcal O_{E_n}[1/7])\{7\},
 \qquad M'=\varprojlim_n C_n',
\]
the kernel of $C_n\to C_n'$ is generated by the seven-prime classes and
is $\Delta$-trivial. It follows that
\[
 (M')^\chi
 \simeq\left(\varprojlim_n C_n\right)^\chi
 \simeq\Lambda/(7,T).
 \tag{S34}\label{seven:split-module}
\]
Class field theory identifies $M'$ with the maximal unramified abelian
pro-seven Galois group over $E_\infty$ in which every seven-prime splits
completely. This is the norm inverse limit of the finite Hilbert
$S$-class-field correspondences
\cite[Remark~1.3.21 and Proposition~3.2.20]{sevenSharifi}.

Here one can check the global and local descent maps explicitly.
Let $G=\Delta\times\Gamma$ and $D=D_{\chi,\kappa}$.
The element $\chi(\delta)-1$ is a unit in $\mathcal O$, so $D^\Delta=0$.
Cohomology in positive degree for the order-three group $\Delta$ vanishes
on seven-primary modules, by averaging.
The Hochschild--Serre sequence therefore gives
$H^i(G,D)=0$ for all $i\ge0$.
The inflation--restriction sequence is consequently an isomorphism
\[
 H^1(K,D)\ \xrightarrow{\ \sim\ }\
 H^1(E_\infty,D)^G
 =\operatorname{Hom}_{\rm cts}(G_{E_\infty},D)^G.
 \tag{S35}\label{seven:global-inflation}
\]
At either seven-prime the decomposition group in $G$ is all of
$\Delta\times\Gamma$: the local extension is the compositum of the
unramified cubic extension and the totally ramified cyclotomic tower.
The identical argument shows that local inflation--restriction is an
isomorphism there. Thus a class vanishes on the full local decomposition
group before descent if and only if its descended homomorphism does.
At a finite prime away from seven the inertia quotient in $G$ has order
one or three; its positive cohomology again vanishes. This makes
restriction injective on the inertia cohomology, so the unramified
condition is also equivalent before and after descent.
These are precisely the two local conditions in
\eqref{seven:strict-definition}. We obtain
\[
 \operatorname{Sel}_{\rm str}(K,D)
 =\operatorname{Hom}_{\rm cts}(M',D)^{\Delta\times\Gamma}.
 \tag{S36}\label{seven:Hom-description}
\]

For clarity about character and inverse conventions, equivariance means
\[
 f(\delta x)=\chi(\delta)f(x),\qquad f(\gamma x)=u f(x).
 \tag{S37}\label{seven:orientation}
\]
The first condition selects the natural $\chi$-component of $M'$.
On that component $\gamma=1$ and $7=0$ by \eqref{seven:split-module}.
Hence \eqref{seven:Hom-description} is $D[7,u-1]$, which, since
$\mathcal O$ is a discrete valuation ring, is noncanonically
$\mathcal O/(7,u-1)$.
This proves \eqref{seven:all-characters}.
Writing a contragredient convention instead replaces $u-1$ by
$u^{-1}-1=-u^{-1}(u-1)$, the same ideal. No Tate twist or unspecified
specialization kernel occurs in this argument.

\subsection{Weights, wild characters, and coefficient normalization}
\label{seven:weights}

For $e\in\mathcal O$, the convergent exponential gives the continuous
character determined by
$\kappa(\gamma)=\exp(-e\log8)$.
It satisfies $u-1\in7\mathcal O$. In particular this includes
$\kappa=\langle\varepsilon_7\rangle^{-e}$ for $e\in\mathbb Z_7$, including
$e=0$, and
\[
 \operatorname{Sel}_{\rm str}
   \bigl(K,(V/\mathcal O)(\chi\langle\varepsilon_7\rangle^{-e})\bigr)
 \simeq\mathcal O/7\mathcal O.
 \tag{S38}\label{seven:principal-weights}
\]
Over $\mathbb Z_7$ every such group has length one.
Over a ramified coefficient ring its $\mathcal O$-length is
$e(V/\mathbb Q_7)$; it must not be replaced by a length-one quotient by
a uniformizer.

More generally, suppose $\mathcal O$ contains a primitive $7^r$th root
$\zeta$, where $r\ge1$, and put
$u=\zeta\exp(-e\log8)$ with $e\in\mathcal O$.
Since
$v_7(\zeta-1)=1/(6\cdot7^{r-1})<1$
and $\exp(-e\log8)-1\in7\mathcal O$, the two summands in
$u-1=(\zeta-1)+\zeta(\exp(-e\log8)-1)$ have different valuations.
Thus
\[
 \operatorname{Sel}_{\rm str}(K,D_{\chi,\kappa})
 \simeq\mathcal O/(\zeta-1),\qquad
 \frac{\log\#\operatorname{Sel}_{\rm str}(K,D_{\chi,\kappa})}
      {[V:\mathbb Q_7]}
 =\frac{\log7}{6\cdot7^{r-1}}.
 \tag{S39}\label{seven:wild-cost}
\]
The valuation formula follows by applying the Eisenstein polynomial
$\Phi_{7^r}(1+X)$ to $\zeta-1$.
For the principal branch \eqref{seven:principal-weights}, the same
normalized logarithmic size is $\log7$.
For an arbitrary continuous $\kappa$, the completely general answer
remains \eqref{seven:all-characters}; no parametrization of all such
characters by $\exp(-e\log8)$ is assumed.

Finally consider a coefficient family written in the original finite
character convention as
\[
 D_{j,a,e}=(V/\mathcal O)
    \bigl(\chi^j\xi_7^a\varepsilon_7^{-e}\bigr),\qquad
 j\in\{1,2\},\quad0\le a<6,\quad e\ge1,
 \qquad \xi_7=\omega_7^v,\quad v\in\{1,5\}.
 \tag{S40}\label{seven:matched-coefficients}
\]
If the finite residual characters are matched, then
$\xi_7^a\omega_7^{-e}=1$ exactly, because reduction is injective on
$\mu_6$. Equivalently, $va\equiv e\pmod6$, which determines one $a$
for every $e$. Thus each matched coefficient is exactly
$(V/\mathcal O)(\chi^j\langle\varepsilon_7\rangle^{-e})$, and
\[
 s(e)=1\quad\hbox{for every positive matched weight},\qquad
 \sum_{j=1}^{2}\sum_{e=1}^{N}
 \frac{\log\#\operatorname{Sel}_{\rm str}(K,D_{j,a(e),e})}
      {[V:\mathbb Q_7]}
 =2N\log7.
 \tag{S41}\label{seven:matched-cost}
\]
There is no exceptional weight in this matched family. Statements about
other tame characters or additional factors require their own local
matching and are not used here.

\section{Universal norms and the linear one-prime factor}
\subsection{Capitulation and norms of global units}

All module notation for completed units in this proof is additive.
Theorem~\ref{seven:main} gives
$A(E_n)^\chi\simeq C_7$ and norm isomorphisms between these groups.
For $n>m$, extension and norm of ideals satisfy
\[
 N_{n/m}j_{n/m}=7^{n-m}
 \quad\hbox{on } A(E_m)^\chi.
\]
Since the right side vanishes and the norm is injective, we have
\begin{equation*}\tag{U3}\label{seven:faithful-capitulation}
 j_{n/m}=0.
\end{equation*}

\begin{lemma}\label{seven:unit-norm-index}
For every $n>m$,
\[
 H^1(\operatorname{Gal}(E_n/E_m),\mathcal E_n)^\chi\simeq C_7,
 \qquad
 [\mathcal E_m^\chi:N_{n/m}\mathcal E_n^\chi]=7.
\]
In particular
$N_{n/0}\mathcal E_n^\chi=7\mathcal E_0^\chi$ for $n\ge1$.
\end{lemma}

\begin{proof}
Set $L=E_n$, $k=E_m$, and $G=\operatorname{Gal}(L/k)$. Hilbert 90 and
the principal-ideal sequence identify
$H^1(G,\mathcal O_L^\times)$ with $P_L^G/P_k$.
The natural map
\[
 P_L^G/P_k\longrightarrow I_L^G/jI_k
 \simeq
 \bigoplus_{\mathfrak q\ {\rm ramified\ in}\ L/k}
           \mathbb Z/e_{\mathfrak q}\mathbb Z
\]
has kernel equal to the capitulation kernel on ideal classes.
The only ramified primes are the two primes above seven, and each is
fixed by $\Delta$. The displayed direct sum is therefore
$\Delta$-trivial. Applying a faithful idempotent and then completing
at seven proves
\[
 H^1(G,\mathcal E_n)^\chi
 \simeq\ker(j_{n/m})^\chi\simeq C_7.
\]
Completion commutes with this cohomology because $\mathbb Z_7$ is flat
over $\mathbb Z$ and the group is finite.

The complex places of $E_n$ form a regular set for
$\Delta\times\operatorname{Gal}(E_n/E)$. Indeed, each complex place has
exactly one embedding extending a fixed embedding of the imaginary
quadratic field $K$. Dirichlet's unit theorem consequently gives the
regular representation with one trivial representation removed.
On a faithful character the removed term disappears, so
\[
 \mathbb Q_7\otimes_{\mathbb Z_7}\mathcal E_n^\chi
 \simeq
 \mathbb Q_7[\operatorname{Gal}(E_n/E)].
\]
Its restriction to $G$ is a sum of regular representations. Hence the
Herbrand quotient of $\mathcal E_n^\chi$ is one: an integral regular
lattice has trivial Tate cohomology, and replacing a lattice by a
commensurable one does not change its Herbrand quotient, since finite
modules have quotient one. Thus
\[
 \#\widehat H^0(G,\mathcal E_n^\chi)
 =\#H^1(G,\mathcal E_n^\chi)=7.
\]
The invariant units are $\mathcal E_m^\chi$, proving the asserted
norm index. At level zero $\mathcal E_0^\chi$ has rank one, so its
submodule of index seven is $7\mathcal E_0^\chi$.
\end{proof}

We record explicitly the base localization. Put
\[
 \varepsilon_0=-270+739\alpha-108\alpha^2.
\]
Reduction using $\alpha^3=269$ gives
\[
 \varepsilon_0^{342}
 \equiv1+7\alpha+28\alpha^2\pmod{49}.
\]
Since $342\equiv6\pmod7$, the principal logarithm therefore satisfies
\begin{equation*}\tag{U4}\label{seven:base-local-generator}
 \frac{\log\varepsilon_0}{7}
 \equiv6\alpha+3\alpha^2\pmod7.
\end{equation*}
The field $L_0$ is the unramified cubic extension of $\mathbb Q_7$,
and logarithm identifies $U_0$ with $7\mathbb Z_7[\alpha]$.
The two faithful lines have bases $7\alpha$ and $7\alpha^2$.
Both coefficients in \eqref{seven:base-local-generator} are units.
As $\mathcal E_0^\chi$ has rank one, this proves
\begin{equation*}\tag{U5}\label{seven:base-local-isomorphism}
 \ell_0:\mathcal E_0^\chi\xrightarrow{\ \sim\ }U_0^\chi.
\end{equation*}

\subsection{Inverse limits and their control maps}

We use two elementary facts about compact modules. First, inverse
limit is exact on inverse systems of compact Hausdorff abelian groups.
To lift a compatible element of a quotient system, any finite set of
compatibility conditions can be met by lifting at its highest level
and taking images at lower levels. Compactness then supplies a lift
meeting all conditions. Transition maps need not be surjective.
In particular, for a compatible endomorphism $T$,
\begin{equation*}\tag{U6}\label{seven:compact-coinvariants}
 (\varprojlim B_n)/T(\varprojlim B_n)
 \simeq\varprojlim(B_n/TB_n).
\end{equation*}
Indeed, exactness applied first to $B_n\to TB_n$ proves
$\varprojlim TB_n=T\varprojlim B_n$, and a second application proves
\eqref{seven:compact-coinvariants}.

Second, the following compact form of Nakayama does not presuppose
finite generation. If a compact pro-seven $\Lambda$-module $B$ has
continuous finite $\Gamma$-stable quotients separating its points and
$B/(7,T)B$ is generated by the image of $b$, then $B=\Lambda b$.
The submodule $\Lambda b$ is compact and closed. Its quotient $Q$
satisfies $Q=(7,T)Q$. Every continuous finite quotient of $Q$ is a
module for some $\mathbb Z_7[C_{7^r}]$ on which $(7,T)$ is nilpotent,
and so must vanish. The finite-level unit groups modulo powers of
seven give the separating quotients needed here.

\begin{lemma}\label{seven:local-unit-control}
For $n>m$, norm induces an isomorphism
\[
 U_n^\chi/(\gamma^{7^m}-1)U_n^\chi
 \xrightarrow{\ \sim\ }U_m^\chi.
\]
Moreover $U_\infty^\chi\simeq\Lambda$, and
$U_\infty^\chi/\omega_mU_\infty^\chi\simeq U_m^\chi$.
\end{lemma}

\begin{proof}
The local extension $L_n/L_m$ is cyclic and totally ramified.
Local class field theory identifies
$U_m/N_{n/m}U_n$ with its Galois group. The $\Delta$-action on that
quotient is trivial because the unramified cubic and cyclotomic
ramified Galois groups commute. Hilbert 90 and the valuation sequence
also give
\[
 H^1(\operatorname{Gal}(L_n/L_m),\mathcal O_{L_n}^\times)
 \simeq\mathbb Z/7^{n-m}\mathbb Z
\]
with trivial $\Delta$-action. Replacing full units by principal units
does not change these Tate groups, since the residue unit group has
order $342$, prime to seven. Both
$\widehat H^0$ and $H^1$ therefore vanish on the faithful component.
The claimed finite norm isomorphism follows from the cyclic Tate
complex.

Compact exactness gives
$U_\infty^\chi/TU_\infty^\chi\simeq U_0^\chi\simeq\mathbb Z_7$.
Compact Nakayama makes $U_\infty^\chi$ cyclic over $\Lambda$.
All the finite local norms on this character are surjective, so a
generator of the limit maps to a generator of every $U_n^\chi$.
These groups are free of rank $7^n$ over $\mathbb Z_7$. There are no
seven-power roots of unity, since the local ramification index is a
power of seven, and their rational rank follows from the local
normal-basis representation and logarithm. Consequently the
surjection
\[
 \Lambda/\omega_n\Lambda\longrightarrow U_n^\chi
\]
is an isomorphism between free $\mathbb Z_7$-modules of the same rank.
The annihilator of a generator of $U_\infty^\chi$ lies in every
$\omega_n\Lambda$. But $\omega_n\in(7,T)^{n+1}$, and their intersection
is zero. This proves freeness. The same finite norm argument and
compact exactness give control at every level $m$.
\end{proof}

The global control map requires the capitulation terms, including
their transition maps.
For $G_n=\operatorname{Gal}(E_n/E)$ and $n\ge1$, the cyclic Tate
complex and Lemma~\ref{seven:unit-norm-index} give
\begin{equation*}\tag{U7}\label{seven:global-coinvariant-sequence}
 0\longrightarrow H^1(G_n,\mathcal E_n)^\chi
 \longrightarrow\mathcal E_n^\chi/T\mathcal E_n^\chi
 \xrightarrow{\ N_{n/0}\ }7\mathcal E_0^\chi
 \longrightarrow0.
\end{equation*}
The left term has order seven, but its norm transition from any higher
positive level to a lower one is zero. To see this, represent a base
class by an ideal $\mathfrak a$ and choose
$(z_n)=\mathfrak a\mathcal O_{E_n}$. Its cohomology class is represented
by $\gamma(z_n)/z_n$. For $n>m$, put $z_m'=N_{n/m}z_n$.
Then
\[
 (z_m')=\mathfrak a^{7^{n-m}}\mathcal O_{E_m}.
\]
The base class has order dividing seven, so
$\mathfrak a^{7^{n-m}}=(b)$ for some $b\in E^\times$.
Thus $z_m'/b$ is a unit, and
\[
 N_{n/m}\!\left(\frac{\gamma(z_n)}{z_n}\right)
 =\frac{\gamma(z_m'/b)}{z_m'/b}
\]
is a coboundary. This calculation commutes with the faithful
idempotent. The transition on the right of
\eqref{seven:global-coinvariant-sequence} is the identity.
Taking compact inverse limits therefore proves
\begin{equation*}\tag{U8}\label{seven:global-limit-coinvariants}
 \mathcal E_\infty^\chi/T\mathcal E_\infty^\chi
 \xrightarrow{\ \sim\ }7\mathcal E_0^\chi.
\end{equation*}
In particular there is no unaccounted kernel in this specialization.
Compact Nakayama now proves that $\mathcal E_\infty^\chi$ is cyclic.

The finite maps $\mathcal E_n\to\overline{\mathcal E}_n$ are surjective
maps of compact groups and commute with norms. Compact exactness
identifies $\overline{\mathcal E}_\infty$ with the image of
$\mathcal E_\infty\to U_\infty$, without requiring injectivity of
finite-level localization. Choose cyclic generators $e_\infty$ and
$u_\infty$ of $\mathcal E_\infty^\chi$ and $U_\infty^\chi$, and write
\[
 \ell_\infty(e_\infty)=f_{\mathfrak p,\chi}(T)u_\infty.
\]
Equations \eqref{seven:base-local-isomorphism} and
\eqref{seven:global-limit-coinvariants} give
\begin{equation*}\tag{U9}\label{seven:unit-factor-augmentation}
 f_{\mathfrak p,\chi}(0)\in7\mathbb Z_7^\times.
\end{equation*}
Thus $f_{\mathfrak p,\chi}\ne0$. Any relation on $e_\infty$ would
annihilate this nonzero element in the free local module $\Lambda$.
As $\Lambda$ is a domain, the global cyclic module is free and its
localization is injective. We have proved
\begin{equation*}\tag{U10}\label{seven:unit-factor}
 \mathcal E_\infty^\chi\simeq\Lambda,\qquad
 V_{\mathfrak p,\chi}\simeq\Lambda/(f_{\mathfrak p,\chi}).
\end{equation*}
The augmentation condition alone allows either $\mu=1,\lambda=0$
or $\mu=0,\lambda\ge1$. A first-layer calculation distinguishes them.

\subsection{An exact first-layer local index}

Use the period $\beta$, the automorphism $\tau$, and the thirteen
units of \S\ref{seven:finite-data}. Put $\pi=\beta+1$. The period
polynomial gives the Eisenstein equation
\begin{equation*}\tag{U11}\label{seven:local-eisenstein}
 \pi^7-7\pi^6+49\pi^4-98\pi^2-49\pi+7=0.
\end{equation*}
Thus $v_\pi(7)=7$, the residue field is
$\mathbb F_7[\overline\alpha]$ with $\overline\alpha^3=3$, and
$7/\pi^7\equiv-1\pmod\pi$. Set $U^j=1+\pi^j\mathcal O_{L_1}$.

\begin{lemma}\label{seven:local-finite-basis}
Let $\chi_a(\delta)=\zeta_3^a$, where
$\delta(\alpha)=\zeta_3\alpha$, $a=1,2$, and
$\zeta_3\equiv2\pmod7$. For
$J=\{1,2,3,4,5,6,8\}$, the seven elements
\[
 b_{a,j}=e_{\chi_a}(1+\alpha^a\pi^j),\qquad j\in J,
\]
form a $\mathbb Z_7$-basis of $U_1^{\chi_a}$.
Their images in $(U^1/U^{16})^{\chi_a}$ have respective orders
$343,49,49,49,49,49,49$ and form a direct-product basis.
\end{lemma}

\begin{proof}
For $i=0,1,2$, the successive leading grades of
$g_i=1+\alpha^i\pi$, $g_i^7$, and $g_i^{49}$ are $1,7,14$,
with coefficients
$\overline\alpha^i,\overline\alpha^{7i},-\overline\alpha^{7i}$.
For $h_{i,j}=1+\alpha^i\pi^j$, with
$j\in\{2,3,4,5,6,8\}$, the leading grades of $h_{i,j}$ and
$h_{i,j}^7$ are $j,j+7$, with coefficients
$\overline\alpha^i,-\overline\alpha^i$.
These assertions follow from the binomial expansion: the competing
lowest grades in $(1+c\pi^j)^7-1$ are $7+j$ and $7j$.
The next powers lie in $U^{16}$.

The displayed pivot grades partition $\{1,\ldots,15\}$, and the three
coefficients at each grade form an $\mathbb F_7$-basis of the residue
field. Successive digit elimination gives
\[
 U^1/U^{16}\simeq(C_{343})^3\oplus(C_{49})^{18}
\]
with these generators. On each associated graded quotient,
$e_{\chi_a}$ projects onto
$\mathbb F_7\overline\alpha^a$.
Moreover
$\overline\alpha^{7a}=3^{2a}\overline\alpha^a$ stays in that line.
The same elimination proves the asserted faithful direct-product
basis.

Logarithm and exponential are inverse on $U^j$ for $j\ge2$, since
$j>7/(7-1)$. Hence
$U^{16}=(U^9)^7\subseteq(U^1)^7$.
Reduction modulo $U^{16}$ therefore preserves
$U_1^{\chi_a}/7U_1^{\chi_a}$. The seven elements generate this quotient,
and $U_1^{\chi_a}$ is free of rank seven over $\mathbb Z_7$.
Nakayama proves that the displayed lifts form a full basis.
\end{proof}

The seven projected global units to compare with this basis are
\[
 e_{\chi_a}\varepsilon_0,\quad
 e_{\chi_a}\eta,\quad e_{\chi_a}\tau\eta,\quad\ldots,\quad
 e_{\chi_a}\tau^5\eta.
\]
They form a full $\mathbb Z_7$-basis of $\mathcal E_1^{\chi_a}$.
Indeed, prime-to-seven descent identifies
$\mathcal E(F_1)=\mathcal E(E_1)^{c=1}$, and projection gives an exact
sequence
\[
 0\longrightarrow\mathcal E(Q_1)
 \longrightarrow\mathcal E(F_1)
 \xrightarrow{\ e_{\chi_a}\ }\mathcal E(E_1)^{\chi_a}
 \longrightarrow0.
\]
For the kernel, the $\Delta$-trivial unit module is
$\mathcal E(KQ_1)$; its units are fixed by $c$ after seven-adic
completion. In fact $u/c(u)$ has absolute value one at all complex
embeddings and is a root of unity, and norm and extension identify
the completed units of this CM field with those of $Q_1$.
For surjectivity a vector $y$ in the faithful line lifts to $y+cy$.
The first six units in the thirteen-element seven-saturated basis
lie in the rank-six kernel. Their span is a direct summand, and hence
saturated, so it is the entire kernel. The remaining seven vectors
therefore give the stated basis.

Raise these seven global units to $342$ to make them principal.
This multiplies their completed coordinate columns by a
$\mathbb Z_7$-unit, and does not change their lattice index.
In the local basis of Lemma~\ref{seven:local-finite-basis}, the
coordinate matrices modulo $49$, with row order
$1,2,3,4,5,6,8$ and column order
$\varepsilon_0,\eta,\tau\eta,\ldots,\tau^5\eta$, are
\begin{equation*}\tag{U12}\label{seven:local-matrix-one}
 B_1=
 \begin{pmatrix}
 21&0&28&14&35&14&35\\
 0&48&48&34&48&27&48\\
 0&17&12&0&2&11&27\\
 28&45&24&25&41&2&34\\
 0&48&39&34&29&48&3\\
 7&26&15&30&16&22&12\\
 3&44&4&40&9&8&0
 \end{pmatrix},
 \qquad\det B_1\equiv28\pmod{49},
\end{equation*}
and
\begin{equation*}\tag{U13}\label{seven:local-matrix-two}
 B_2=
 \begin{pmatrix}
 42&42&21&0&14&21&28\\
 0&30&23&9&9&16&37\\
 0&20&44&33&29&32&28\\
 21&37&43&19&42&14&33\\
 0&22&32&31&13&14&28\\
 28&11&29&0&35&23&37\\
 13&44&0&12&34&23&32
 \end{pmatrix},
 \qquad\det B_2\equiv7\pmod{49}.
\end{equation*}

For precision, the finite arithmetic needed to verify these entries
can be performed entirely in successive $\pi$-adic digits. Equation
\eqref{seven:local-eisenstein} gives the exact carry rule
\begin{equation*}\tag{U14}\label{seven:local-carry}
 7\equiv-\pi^7-\pi^{13}+\pi^{15}\pmod{\pi^{16}}.
\end{equation*}
Use sixteen digits, each in the basis
$1,\alpha,\alpha^2$ with coefficients in $\{0,\ldots,6\}$.
Reduce $\alpha^3=269$ and carry every coefficient multiple of seven
upward by \eqref{seven:local-carry}; each carry strictly increases its
$\pi$-degree. This gives the ring
$\mathcal O_{L_1}/(\pi^{16})$ of order $7^{48}$.
All denominators in the printed unit expressions are prime to seven,
and their inverses modulo $343$ suffice, since $7^3=0$ in this ring.
Use $\beta=\pi-1$, the printed polynomial for $\tau$, and
$\zeta_3=324\pmod{343}$. On a principal unit the projector is the
product of its three $\Delta$-conjugates with exponents
$3^{-1}\zeta_3^{-ar}\pmod{343}$, $r=0,1,2$.
The pivot grades in the preceding lemma then recover the first
coordinate modulo $343$ and the other six modulo $49$.
This procedure determines all entries of
\eqref{seven:local-matrix-one}--\eqref{seven:local-matrix-two} from the
unit data already printed in \S\ref{seven:finite-data}.

Both determinants have valuation exactly one. The local and global
vectors used are full $\mathbb Z_7$-bases, so we have proved the exact
indices
\begin{equation*}\tag{U15}\label{seven:first-layer-local-index}
 [U_1^\chi:\overline{\mathcal E}_1^\chi]=7
\end{equation*}
for both faithful characters. The calculation applies at either
chosen prime above seven; changing the image of $\omega$ interchanges
the character labels.
As a further exact norm check on the same data, the printed expression
$\eta(\alpha,\beta)$ satisfies
\[
 \operatorname{Res}_B(q(B),\eta(\alpha,B))=1
 \quad\hbox{in }\mathbb Q[\alpha]/(\alpha^3-269).
\]
The six $Q_1$-units have relative norms
$-1,-1,1,-1,-1,1$, and $\varepsilon_0$ has relative norm
$\varepsilon_0^7$. This confirms the first-layer global norm defect
without serving as an additional premise for the all-level argument.

\subsection{The linear factor and the weight valuations}

For every $n>1$, Lemma~\ref{seven:unit-norm-index} gives
$[\mathcal E_1^\chi:N_{n/1}\mathcal E_n^\chi]=7$.
These norm images form a descending chain with the same index, hence
are equal. Compactness identifies the common image with the
projection of $\mathcal E_\infty^\chi$ to level one.
Lemma~\ref{seven:local-unit-control} therefore gives
\[
 V_{\mathfrak p,\chi}/\omega_1V_{\mathfrak p,\chi}
 \simeq
 U_1^\chi/
 \ell_1\bigl(\operatorname{im}
       (\mathcal E_\infty^\chi\to\mathcal E_1^\chi)\bigr).
\]
Equation \eqref{seven:first-layer-local-index} also proves that
$\ell_1$ is injective on this character, since its source and target
have equal rank. The universal-norm image has index seven in the
global units, so
\begin{equation*}\tag{U16}\label{seven:unit-factor-first-quotient}
 \#\Lambda/(f_{\mathfrak p,\chi},\omega_1)=7^2.
\end{equation*}

We now determine the Weierstrass invariants. If all coefficients of
$f=f_{\mathfrak p,\chi}$ are divisible by seven, then
\eqref{seven:unit-factor-augmentation} implies $f=7u$ for a unit $u$.
The quotient in \eqref{seven:unit-factor-first-quotient} would then
have order $7^7$, a contradiction. Otherwise Weierstrass preparation
writes $f=ug$, where $g$ is a distinguished polynomial of degree
$d\ge1$. Its constant coefficient has valuation one, so $g$ is
Eisenstein. If $1\le d\le5$, then at a primitive seventh root
$\zeta$ the leading term of $g(\zeta-1)$ has uniquely smallest
valuation $d/6$. Including the root $\zeta=1$, we obtain
\[
 v_7\!\left(\prod_{\zeta^7=1}g(\zeta-1)\right)=1+d.
\]
If $d\ge6$, each nontrivial-root factor has valuation at least one,
so the displayed valuation is at least seven, or the quotient is
infinite. Equation \eqref{seven:unit-factor-first-quotient} therefore
forces $d=1$. We have proved
\[
 f_{\mathfrak p,\chi}(T)
 =u(T)(T-a_{\mathfrak p,\chi}),\qquad
 a_{\mathfrak p,\chi}\in7\mathbb Z_7^\times.
\]

Rubin's class-field sequence
\cite[printed p.~36, equation~(5)]{Rubin} gives
\eqref{seven:one-prime-exact}. Its quotient is the finite faithful
class module of Theorem~\ref{seven:main}. Thus
$X_{\mathfrak p,\infty}^\chi$ is a finitely generated
$\mathbb Z_7$-module of rank one.
By \cite[Theorem~5.3(v), printed pp.~40--41]{Rubin},
$X_{\mathfrak p,\infty}$ has no nonzero pseudo-null submodule;
in the one-variable setting this excludes every finite submodule.
The $\mathbb Z_7$-torsion of $X_{\mathfrak p,\infty}^\chi$ is finite
and $\Lambda$-stable, so it vanishes. Hence this module is
$\mathbb Z_7$-free of rank one. Its $\Gamma$-action agrees, on the
finite-index submodule $V_{\mathfrak p,\chi}$, with multiplication
by $1+a_{\mathfrak p,\chi}$, and must agree on the whole module.
The submodule of index seven is consequently
$7X_{\mathfrak p,\infty}^\chi$. This proves all the asserted module
identifications.

The hypotheses of Rubin's theorem are satisfied: $K$ has class
number one and six roots of unity, $E/K$ is abelian of degree three,
and $E_\infty/K$ is abelian with
$\operatorname{Gal}(E_\infty/E)\simeq\mathbb Z_7$.
The statement concerns any such one-variable extension, so it
applies to the cyclotomic line; it is not restricted to the
$\mathfrak p$-ramified line.

It remains to evaluate the linear factor. Both
$\log(1+a_{\mathfrak p,\chi})$ and $\log8$ have valuation one, so
$e_{\mathfrak p,\chi}^{*}$ is a seven-adic unit and
$8^{-e_{\mathfrak p,\chi}^{*}}=1+a_{\mathfrak p,\chi}$.
For every $e\in\mathbb Z_7$,
\[
 8^{-e}-(1+a_{\mathfrak p,\chi})
 =(1+a_{\mathfrak p,\chi})
   \left(8^{-(e-e_{\mathfrak p,\chi}^{*})}-1\right).
\]
The logarithm and exponential on $7\mathbb Z_7$ give
$v_7(8^z-1)=1+v_7(z)$ for $z\in\mathbb Z_7$.
A unit in $\Lambda$ evaluates to a unit at every point of
$7\mathbb Z_7$. Equation
\eqref{seven:one-prime-weight-valuation} follows for every
characteristic generator.
Positive integers approximate any element of $\mathbb Z_7$ to
arbitrary precision, and hence these valuations are unbounded.

The standard Iwasawa involution replaces the root
$a_{\mathfrak p,\chi}$ by
$(1+a_{\mathfrak p,\chi})^{-1}-1$, which still has valuation one,
and replaces $e_{\mathfrak p,\chi}^{*}$ by its negative.
It does not remove the discrepancy with the constant strict length.
Thus identifying the strict length with the valuation of a generator
of this one-prime characteristic ideal is false, independently of
unit ambiguity and this involution.
An identification with a particular normalized Katz function still
requires its precise character, period, reciprocity, and Euler-factor
conventions. The conclusion here concerns the characteristic
generator of the explicitly defined one-prime module.

\section{The exact contribution to an all-prime arithmetic sum}
The strict Selmer formula determines the complete cost of the matched seven-adic family. We prove the exact comparison in Corollary~\ref{cor:seven-budget}, whose hypotheses and numerical slack are stated in Section~1.


\begin{proof}[Proof of Corollary~\ref{cor:seven-budget}]
For each of the two labels and every weight, the strict group is $\mathcal O/7\mathcal O$. Its cardinality is $7^b$, including when $\mathcal O$ is ramified over $\Z_7$. Each normalized term is therefore
$b^{-1}\log(7^b)=\log7$. Adding the $2N$ distinct terms proves (B1). Leaving one pair proves (B2). The triangle inequality proves the non-strict inequality in (B3). Since $7<2^3$, one has $2\log7<6\log2$, so
\[
2(T-1)\log7<6T\log2\le6\mathcal P.
\]
This gives its strict inequality. Subtracting the known summand in (B1) proves (B4).
\end{proof}

For example, the arithmetic budget associated with an abelian field $F$ of degree $D$, conductor dividing $f$, discriminant of absolute value $\Delta$, and localization integer $R$ is
\[
 \mathcal P(F,T)=DT^2\log(2DTfR)+TD^4\log(2D\Delta R).
\]
All parameters are positive integers. Thus $\mathcal P(F,T)\ge T\log2$, which is the only property used in the numerical constant in (B3). The two matched labels have trivial auxiliary quadratic and two-power factors, so an application removes exactly those labels; every other term stays in $\mathcal R_{\rm rest}$.

Here is the full effect on a late-window premise. Fix $0<\varepsilon<1/256$, put $T(u)=\lfloor u^{399}\rfloor$, and suppose the existing comparison error $E(u)+C\mathcal P(u)$ is bounded, uniformly in all the prescribed finite data, by
\[
 O_\varepsilon\!\left(u^{806+16\varepsilon}(\log u)^2\right).
\]
The additional explicit summand in (B3) has the same envelope. If constants $J\ge1$, $X_0$ and $s\in(0,1)$ give, for every prescribed pair at each $X\ge X_0$, a window
\[
 X^{1/J}\le W,\quad W^2\le X,\qquad
 \mathcal R_{\rm rest}(u)\le(1-s)\exp(u^\varepsilon/2)
 \quad(W\le u\le W^2),
\]
then the comparison gives $I(u)\le(1-s/2)\exp(u^\varepsilon/2)$ throughout that same window for all sufficiently large $X$. Indeed the polynomial envelope divided by $\exp(u^\varepsilon/2)$ tends uniformly to zero; choose $U_0$ so that it is at most $s/2$ for $u\ge U_0$, and increase $X_0$ to at least $U_0^J$. No new window is selected. The exact finite-threshold adjustment remains (B4).

The largest matched weight therefore ceases to be an unknown arithmetic input. The remaining assertion is a bound for $\mathcal R_{\rm rest}$ with the specified uniform window quantifiers, together with the geometric realization hypotheses of the original comparison. Neither the class-group stability theorem nor (B1) supplies that bound, and the cubic Riemann hypothesis remains open.


\section{What nested exceptional classes alone would imply}
Theorem~\ref{seven:main} shows that the second and every higher length superlevel is empty for this arithmetic family. The following elementary results isolate the weaker conclusion available when only nestedness is known. They explain exactly why the finite class calculation is needed, and are not assertions that exceptional weights occur here.

Let $p$ be a prime and let $s(e)$ be a finite positive integer for each positive integer $e$. Suppose each higher superlevel is either empty or one full residue class modulo $p^{n-1}$. Nested classes either stop, giving a finite ceiling, or converge to one $p$-adic integer. Reading these superlevels gives the entire length profile. Long gaps between the nonzero digits of that integer can produce very large isolated peaks; nestedness by itself excludes none of them.

\begin{theorem}\label{th:profile}
Adopt $v_p(0)=+\infty$. Let $p$ be a prime and $s:\Z_{\ge1}\to\Z_{\ge1}$ a finite-valued function. Suppose that for each $n\ge2$, the set $\{e:s(e)\ge n\}$ is empty or one full residue class modulo $p^{n-1}$. Then precisely one of the following descriptions applies.
\begin{enumerate}
\item The lengths are bounded, with maximum $R\ge1$. There exists $\alpha\in\Z_p$ such that
\[
 s(e)=\min\{R,1+v_p(e-\alpha)\}\qquad(e\ge1).
 \tag{2}
\]
For $R\ge2$, only the residue of $\alpha$ modulo $p^{R-1}$ is specified.
\item The lengths are unbounded. There is a unique $\alpha\in\Z_p\setminus\Z_{\ge1}$ such that
\[
 s(e)=1+v_p(e-\alpha)\qquad(e\ge1).
 \tag{3}
\]
\end{enumerate}
If also $s(e)=1$ whenever $p\mid e$, then the parameter in the unbounded case is a unit, and in the bounded case with $R\ge2$ its specified residue is a unit.
\end{theorem}
\begin{proof}
Superlevel sets are nested by definition. If one is empty, all subsequent ones are empty. Let $R$ be the largest index of a nonempty superlevel set, interpreting $\{s\ge1\}$ as the whole set. If $R=1$, any $\alpha$ gives (2). Otherwise choose $\alpha$ in the last residue class modulo $p^{R-1}$. Nestedness identifies each earlier class with its reduction. Thus for $2\le n\le R$, the inequality $s(e)\ge n$ is equivalent to $v_p(e-\alpha)\ge n-1$; there are no later superlevels. This is (2).

If no superlevel is empty, the compatible residues have a unique inverse limit $\alpha\in\Z_p$. For every $e\ge1$ and every $n\ge2$, membership in the $n$th superlevel is exactly the congruence $e\equiv\alpha\pmod{p^{n-1}}$. Finiteness of $s(e)$ excludes $e=\alpha$ and proves (3). Conversely, (3) determines all residues and hence determines $\alpha$ uniquely. The assertion about units follows already from the second superlevel.
\end{proof}

\begin{proposition}\label{prop:peaks}
For $p=7$, the hypotheses of Theorem~\ref{th:profile}, including $s(e)=1$ when $7\mid e$, do not imply a polynomial upper bound for $\max_{1\le e\le N}s(e)$.
More precisely, there is such a function for which, for every real $A>0$,
\[
 \limsup_{N\to\infty}\frac{\max_{1\le e\le N}s(e)}{N^A}=\infty.
 \tag{4}
\]
This is a counterexample to an implication between abstract length conditions, not a claimed arithmetic Selmer example.
\end{proposition}
\begin{proof}
Set $N_1=1$, $N_{k+1}=7^{N_k^2}$, and
\[
 \alpha=1+\sum_{k\ge1}7^{N_k}\in\Z_7.
\]
The series converges seven-adically, defines a unit, and is not a positive integer, because its base-seven expansion has infinitely many nonzero digits. Define $s$ by (3). It satisfies all the stated superlevel and divisibility conditions. At the positive integer
$e_k=1+\sum_{j\le k}7^{N_j}$ one has
\[
 e_k<2\cdot7^{N_k},\qquad s(e_k)=1+N_{k+1}=1+7^{N_k^2}.
\]
For fixed $A$, the ratio is at least
$7^{N_k^2-A N_k}/2^A$, which tends to infinity. This proves (4).
\end{proof}

\begin{proposition}[the sum after deleting one maximum]\label{prop:sum}
Under the hypotheses of Theorem~\ref{th:profile}, choose $e_*$ maximizing $s(e)$ for $1\le e\le N$. Then
\[
 \sum_{\substack{1\le e\le N\\e\ne e_*}}s(e)
 \le (N-1)+v_p((N-1)!).
 \tag{5}
\]
\end{proposition}
\begin{proof}
The first layer contributes $N-1$. At a higher layer, either the superlevel does not meet the interval, or it contains $e_*$, by maximality. A residue class modulo $p^k$ has at most $1+\lfloor(N-1)/p^k\rfloor$ points in an interval of $N$ consecutive integers. Deleting $e_*$ therefore leaves at most $\lfloor(N-1)/p^k\rfloor$. Summing over $k\ge1$ and applying Legendre's formula proves (5).
\end{proof}
The two conjugate cubic characters have the same length, so at $p=7$ their combined contribution after this deletion is at most
$2\log7\bigl((N-1)+v_7((N-1)!)\bigr)\le(7/3)(N-1)\log7$.
Proposition~\ref{prop:peaks} explains why the abstract residue conditions alone did not settle the largest length. Theorem~\ref{seven:main} settles it for the present arithmetic family: all its lengths are one.

\section{What a main-conjecture comparison has to identify}
Write $7=\mathfrak p\overline{\mathfrak p}$ in $K$.
For any discrete coefficient $D$, let $V(D)$ be its global first cohomology subject to the fixed inertia conditions away seven. Define
\[
 S_{\mathfrak p}(D)=\ker\bigl(V(D)\to H^1(K_{\mathfrak p},D)\bigr),
 \quad
 S_{\rm str}(D)=\ker\bigl(V(D)\to H^1(K_{\mathfrak p},D)\oplus
 H^1(K_{\overline{\mathfrak p}},D)\bigr).
\]
There is an exact sequence
\[
 0\longrightarrow S_{\rm str}(D)\longrightarrow S_{\mathfrak p}(D)
 \longrightarrow\operatorname{im}\left(S_{\mathfrak p}(D)
 \to H^1(K_{\overline{\mathfrak p}},D)\right)\longrightarrow0.
 \tag{6}
\]
It is simply the kernel-image sequence for the second restriction map.
Whenever the groups are finite, their lengths consequently satisfy
\[
 \ell(S_{\rm str}(D))=\ell(S_{\mathfrak p}(D))-\ell(\operatorname{im}\operatorname{res}_{\overline{\mathfrak p}}).
 \tag{7a}
\]
Thus identifying a one-prime Selmer characteristic ideal with a Katz function does not identify a strict two-prime length until the last term and the specialization errors have been determined.

There is a separate algebraic specialization issue. Put $R=\Z_7[[T]]$ and $C=R/(7,T)$. The modules $0$ and $C$ have the same characteristic ideal, namely the unit ideal, because $C$ has support of codimension two. Yet for every $t\in7\Z_7$, their quotients by $T-t$ have lengths zero and one respectively. Characteristic-ideal equality therefore does not by itself determine the exact order of a finite specialization. One must either eliminate these finite errors, prove a control theorem that cancels them, or work with an appropriate integral Fitting or determinant statement.

Rubin's precise theorem applies more broadly than a one-prime-tower formulation alone suggests. In Rubin (1991, Section 4, pp.~36--37), the finite extension $K_0/K$ is abelian of degree prime to $p$, contains the Hilbert class field, and $K_\infty/K_0$ may be any $\Z_p$-extension or the $\Z_p^2$-extension abelian over $K$. For $K=\Q(\omega)$, $K_0=\Q(\omega,\sqrt[3]{269})$ and $p=7$, all hypotheses hold: $h_K=1$, $7\nmid6\cdot3$, and 7 splits in $K$. In particular his Theorem 4.1(i) applies to the cyclotomic tower. It gives
\[
 \operatorname{char}(A_\infty^\chi)
 =\operatorname{char}((\overline{\mathcal E}_\infty/
                         \overline{\mathcal C}_\infty)^\chi),
\]
where $A_\infty$ is the inverse limit of ideal class groups, $\overline{\mathcal E}_\infty$ denotes global units, and $\overline{\mathcal C}_\infty$ elliptic units.

The distinction relevant to a Katz function is in Rubin's exact sequence (5), p.~36:
\[
0\longrightarrow\overline{\mathcal E}_\infty/
                 \overline{\mathcal C}_\infty
\longrightarrow U_{\mathfrak p,\infty}/
                 \overline{\mathcal C}_\infty
\longrightarrow X_{\mathfrak p,\infty}
\longrightarrow A_\infty\longrightarrow0.
\]
Here $U_{\mathfrak p,\infty}$ uses local units above one chosen prime $\mathfrak p\mid7$, and $X_{\mathfrak p,\infty}$ permits ramification only there. The associated exact sequence after removing elliptic units is
\[
0\longrightarrow U_{\mathfrak p,\infty}/
                  \overline{\mathcal E}_\infty
\longrightarrow X_{\mathfrak p,\infty}
\longrightarrow A_\infty\longrightarrow0.
\tag{7b}
\]
Thus characteristic ideals satisfy
\[
 \operatorname{char}(A_\infty^\chi)=
 \frac{\operatorname{char}(X_{\mathfrak p,\infty}^\chi)}
 {\operatorname{char}((U_{\mathfrak p,\infty}/
                       \overline{\mathcal E}_\infty)^\chi)}.
\tag{8}
\]
Rubin's Remark (1) after Theorem 4.1 identifies the semilocal-unit quotient with a $p$-adic $L$-function in the one-prime and two-variable settings, citing Coates--Wiles and Yager. Vigui\'e (2011, Theorem 1.1, p.~3) gives a separate formulation in the one-prime tower. In the cyclotomic tower considered here, Theorem~\ref{seven:one-prime-linear} computes both the numerator and denominator of (8) as the ideal $(T-a_{\mathfrak p,\chi})$. Their quotient is the unit ideal, consistently with the finite class module.

More strongly, the formula equating $s(e)$ with the valuation of any generator of the one-prime characteristic ideal is false: the former is constantly one and the latter has the unbounded profile \eqref{seven:one-prime-weight-valuation}. This conclusion does not identify a particular Katz measure. Its periods, reciprocity convention, Euler factors, and restriction to the cyclotomic line must still be specified before one can assert a literal formula involving $L_{\rm Katz}$. The unit quotient and the finite specialization correction have now both been determined algebraically.

\section{Finite-level and universal localization}
\begin{proposition}\label{prop:seven-units}
For either prime $\mathfrak p\mid7$ of $K$ and either faithful cubic character $\chi$, localization identifies the completed global unit line of $E$ with the corresponding principal local unit line. Consequently the maximal abelian pro-seven Galois group over $E$ unramified outside the chosen prime has faithful component $C_7$.
\end{proposition}
\begin{proof}
The localization isomorphism is \eqref{seven:base-local-isomorphism}. The finite ray class sequence has kernel the quotient of principal local units by global units and cokernel $A(E)^\chi$. The kernel vanishes on this character, while Theorem~\ref{seven:main} gives $A(E)^\chi=C_7$.
\end{proof}
The finite localization isomorphism does not persist for universal norms: Theorem~\ref{seven:one-prime-linear} computes the nonzero quotient as $\Lambda/(T-a_{\mathfrak p,\chi})$. Thus the denominator in (8) is now determined, including its degree and constant-term valuation.

\section{The exact specialization correction}
This section gives an algebraic formula that must accompany any characteristic-series calculation. Write $\Lambda=\Z_7[[T]]$.

\begin{lemma}\label{lem:euler-specialization}
Let $M$ be a finitely generated torsion $\Lambda$-module and let $f_M$ generate its characteristic ideal. If $t\in7\Z_7$ and $f_M(t)\ne0$, then $M/(T-t)M$ and $M[T-t]$ are finite and
\[
 \operatorname{length}_{\Z_7}M/(T-t)M
 =v_7(f_M(t))+\operatorname{length}_{\Z_7}M[T-t].
\tag{11}
\]
The last group is contained in the maximal finite $\Lambda$-submodule of $M$.
\end{lemma}
\begin{proof}
First take $M=\Lambda/(f)$. Since $T-t$ does not divide $f$, it is a nonzero divisor on this cyclic module. Its cokernel is $\Z_7/(f(t))$, so (11) follows. Direct sums preserve the difference between the two lengths. The structure theorem supplies a pseudo-isomorphism from a direct sum of such cyclic modules to $M$, with finite kernel and cokernel. For a finite module $C$, the kernel and cokernel of any endomorphism have the same order; hence multiplication by $T-t$ has Euler characteristic zero on $C$. The long exact sequence of kernels and cokernels therefore preserves the length difference under the pseudo-isomorphism, proving (11). Finiteness follows at the same time. Finally $M[T-t]$ is finite and stable under $T$ and $\Z_7$, so it is a finite $\Lambda$-submodule. The sum of all finite submodules is finite because $M$ is noetherian; it contains this kernel.
\end{proof}

For the matched character family, take $\Gamma$ to be the cyclotomic seven-extension of $L$, choose $\gamma$ with $\langle\varepsilon_7\rangle(\gamma)=8$, and put $t_e=8^{-e}-1$. Let $A_\infty^\chi$ be the corresponding cubic part of the unramified class module. Imposing complete splitting at primes above 7 does not alter this component. Indeed there is one prime over each prime of $K$ above 7, and the cubic group acts trivially on its unramified local decomposition group; the faithful character projector kills these decomposition groups. Inflation--restriction identifies the strict group with the dual of the specialization at $T=t_e$. Its possible global and local inflation terms vanish on the faithful character component: the cubic character has no invariants globally or at either 7-adic decomposition group. Thus (11), with the previously established finiteness of the strict group, yields the exact equality
\[
 s(e)=v_7(f_A(t_e))+
       \operatorname{length}_{\Z_7}A_\infty^\chi[T-t_e].
\tag{12}
\]
The two faithful characters give the same length by complex conjugation.

If a normalized analytic specialization $\mathcal L_\chi$ has been identified with a generator of $\operatorname{char}(X_{\mathfrak p,\infty}^\chi)$, and $R_\chi$ generates the denominator of (8), the resulting formula is
\[
 s(e)=v_7(\mathcal L_\chi(t_e))-v_7(R_\chi(t_e))+
       \operatorname{length}_{\Z_7}A_\infty^\chi[T-t_e],
\tag{13}
\]
whenever the separate evaluations are nonzero. Formula (13) is conditional only on the specified analytic identification; (11) and (12) are algebraic statements. In the present tower Theorem~\ref{seven:main} gives the stronger integral identification $A_\infty^\chi=\Lambda/(7,T)$. Its characteristic ideal is the unit ideal, so one may take $f_A=1$. Since $t_e\in7\Z_7$, the kernel $A_\infty^\chi[T-t_e]$ is the entire group $C_7$. Thus (12) reads $s(e)=0+1$. The nonzero strict length comes entirely from the finite submodule, which a characteristic ideal cannot detect. There is no exceptional residue for the strict group. Theorem~\ref{seven:one-prime-linear} shows that, whenever $t_e\ne a_{\mathfrak p,\chi}$, the two characteristic valuations in (13) are equal and cancel. At $t_e=a_{\mathfrak p,\chi}$ each separate valuation is infinite, so their difference is not defined; the direct finite-module calculation still gives $s(e)=1$. Formula (13) applies only away from those separate zeros and remains conditional on the stated analytic normalization.

\section*{Method and availability}
The mathematics, text, and original computations were produced by
GPT-6 Astra in ChatGPT Pro mode. Fukuda's stability theorem, Friedman's
regulator bound, the cited class-number and theta identities, and Rubin's
class-field sequence and finite-submodule theorem are named inputs.
The finite witnesses, norm and control arguments, local-unit matrices,
and exact strict and one-prime specialization formulas are supplied here.
The accompanying records distinguish statement-only reconstruction of
class-group stability, exact finite computations, and targeted checks
of the additional module arguments. No human review or formal
verification is claimed. No literal Katz-measure normalization or
bound for the remaining all-prime arithmetic sum is asserted.
The original text and programs are dedicated to the public domain under
CC0 1.0; cited works retain their own terms. Sources and verification
records are at \url{\workbenchurl}, with version records at
\url{\recordurl}.

\appendix
\section{Complete finite unit and residue data}
\label{seven:finite-data}

The following arrays are exact integers. Coefficients for $p$, $q$, $P$ and
$Q$ are in descending order. The coefficients for $\tau$, the six small
units, and the cubic unit are in ascending order. A small-unit numerator
$[a_0,\ldots,a_6]$ represents $\sum a_i\beta^i$; the cubic-unit numerator
represents a polynomial in $\alpha$.
The twenty-one coefficients of $\eta$ are ordered with $\alpha$ varying
first: coefficient number $3j+i+1$ multiplies $\alpha^i\beta^j$,
$0\le i\le2$, $0\le j\le6$.
To keep every integer on the page, each coefficient of $\eta$ is printed as
$[s,a_0,a_1,a_2]$ in base $B_*=10^{100}$, meaning
$s(a_0+B_*a_1+B_*^2a_2)$; no leading zeroes are required.

Each residue triple is $(\ell,r,z)$: reduce $\theta$ to $r$ modulo $\ell$
and use $z$ in \eqref{seven:residue-character}. The thirteen primes with
these residues are distinct, even when their rational primes coincide.
The roots of $p$ are simple at all these residues; the derivative residues,
in the same order, are
\[
 153,14,99,174,74,74,436,46,119,368,219,76,606.
\]
Hensel's lemma therefore supplies the corresponding maps of the full
integer ring to $\mathbb F_\ell$.
The rows of the final matrix are in this order and the columns are
$u_1,\ldots,u_{13}$.
Polynomial substitution, the norm and integrality tests, the finite-field
powers in \eqref{seven:residue-character}, and Gaussian elimination modulo
seven are a complete exact verification of these finite witnesses.

\begingroup
\fontsize{7.5}{9.5}\selectfont
\begin{verbatim}
p = [1, 0, -63, -1820, 1596, 47859, 1550766, -771102, 74382, -773205433, 265573350, -12672069060,
  225337203987, -164087914221, 5350549479537, -37293576067279, 55623366413664, -844644350170257,
  3089185218698537, -5358907459195815, 42933560367107697, -74554880595863287]
q = [1, 0, -21, -21, 91, 112, -84, -97]
P = [7, 0, -252, 43540, 3115, -767592, 28694829, 4854766, -154925106, 3054427187, 508674397, 875124348,
  32367791232, 6622577365]
Q = [28, 0, -819, 72163, 8253, -926394, 25924017, 4555698, -71871954, 1423159430, 242638326,
  1061561823, 4574293518]
tau_numerator = [3210, -2025, -2700, 820, 495, -51, -15]
tau_denominator = 589
small_unit_numerators = [
  [48, 18, -65, -13, 17, 2, -1],
  [-842, 2689, 5156, -305, -1495, -78, 81],
  [-5185, -1631, 4697, 659, -1028, -69, 49],
  [5927, -5217, -7545, 1048, 1825, 44, -91],
  [5085, -3117, -2389, 743, 330, -34, -10],
  [10258, -972, -10131, -431, 2358, 164, -125],
]
small_unit_denominators = [1, 589, 589, 589, 589, 589]
cubic_unit_numerator = [-270, 739, -108]
eta_denominator = 1767
eta_coefficients_in_base_10_power_100 = [
  [1,
   8110910039399052799388839509913995719324441857259887672553997475894562301601700471535299521506405517,
   3431137402558321943275234065603513747732585590776418713813607417676553074824870235139769709481053259,
   3940696968887],
  [1,
   2863301462649584998067112611608264044687998196095409917598114508277713738730199925757154728888953055,
   970812101337790668521585581275454765494152664299295398136805535577832483461611290913002007827311071,
   610457752147],
  [1,
   4821724478914456225350364309117571242884298504341022610343187741267858848133670184073817508713496298,
   3511990174213849653115720097913645282815157493086558732902294533502899394979501897922579875243606017,
   94566689623],
  [1,
   760344159491231074673603899303458724681328958152289201635437873131003916558596294478134311589863496,
   2317480636152653863897711154536450131210677689030596417135280558908466878562464587662862833652562118,
   7059069004615],
  [1,
   5713646533666935748952660526437906230094439209439484821582261626527132028034044010356200033729959598,
   8155918751592805126188738962782659016016516467436799064429460997991815680374480211693583033342107571,
   1093528233921],
  [1,
   5458376416976996449886152263319535978441113085075418487179464286576609680040631434748810087394632375,
   2815314908475193445673806280168229698608681052505411231685686073970229478053287079417408542790496382,
   169399675453],
  [1,
   9104803622062202912042424298875099840356676073442457598635399038039260264973760983790173980634551253,
   5209978760288616211810004718597728762354940779407360157403598658202633153414464523403435572630242052,
   1981995488635],
  [1,
   3229190277068481099567369991097155293895283837272933560810654867460601672084248677453637066047248351,
   1066602590408197039412448163000624315584233985107205159080701191411246693038208294958734638546267908,
   307033126452],
  [1,
   4719833413501908223670674319560071208620416677277259622958912093004163789224972446770714978198517369,
   3589225343881354158630968166458173644607882286774087545514248915094188084702046502205303007048245172,
   47562843245],
  [-1,
   7252159511744921579747663915866179139747715077640636824080064020836768547568447722702965774255379864,
   2117795437844894019014366078716012207469677439109909489100232555667366018470798129278480577073347195,
   1862911672364],
  [-1,
   4596952397808038061504443355452507844075138976341878918188397033987218237849921882684029382396888813,
   5548014240268422805702685958606065400771312345386887361270938503247549652804788711898670621959818004,
   288585719972],
  [-1,
   8576573868236387713297262546743178973121352627702349510026823016588229660225220810296942340315536894,
   7246874522795187239999206975901393322813195843039386133915845845849949086654946548202311744607548915,
   44705134981],
  [-1,
   7281143800513613602674102889822350672969769041691650309307424554802016538032032999848411121384589417,
   7247656128371841488568688653380385878443917715304128569138890771204319093551264064535444911867233869,
   870696576112],
  [-1,
   2603158096105930003776498071993161225442241308288062551734654529270377694879518302937682334035359418,
   8801224015145560570571877684554110083953909461438848969946722322537320157416843524148670087474475623,
   134880575403],
  [-1,
   2144340361888752919586412430689254661717183112022793764572101308645944000985194209283484734442009651,
   4344877732372528348076072331962740358969404598366277209320255591698654981676044669701093523564674725,
   20894500013],
  [1,
   1658789263126353124561691351357597966888983262121120607405319802389834239184987248275993385042870297,
   4563350942981641714892768486786517151934762869785113709687288930688939203417092386426646047883401037,
   47445317587],
  [1,
   3334821144767026713097794436460203379080665880747592876486755438276015329827178694499770466419371729,
   4734315972881480902169887703698274649856456832930081114001007355628760067288435055252831111892960802,
   7349806938],
  [1,
   2179355790946349106219951244546405670330306327971639576268515952647904221008259146456375776668744595,
   7511295224937829629681218123710093110445808535628505578790220369499924491111462384032633379660007227,
   1138566770],
  [1,
   1940103518950861831403240058854220493243755110925587464386601540881397210941906392843448161215669392,
   7398981907903882144917187999221062721013600253004142759479698501891917306885982478145740682084331734,
   43903316832],
  [1,
   3471460889921431092203467331457652901571380143073529228472752287967796312304491085909952970479760717,
   4781793234571883106444295805054726784965688101009565012732394759877052590133372264885207805147053091,
   6801111660],
  [1,
   1294025870817364042382316796528664510562816139774755081398276045466303764150138262507798612436708108,
   7200960424005516814554303735541605542146140398242153380225569364053345312160643463981887310532844365,
   1053567774],
]
prime_root_zeta_triples = [
  [197, 164, 104],
  [197, 119, 104],
  [197, 110, 104],
  [197, 69, 104],
  [197, 50, 104],
  [197, 29, 104],
  [491, 435, 138],
  [491, 431, 138],
  [491, 411, 138],
  [491, 107, 138],
  [491, 98, 138],
  [491, 93, 138],
  [883, 874, 707],
]
residue_character_matrix = [
  [2, 3, 4, 2, 6, 5, 0, 5, 5, 2, 3, 3, 5],
  [5, 6, 2, 5, 1, 1, 0, 5, 5, 5, 2, 3, 3],
  [0, 1, 3, 1, 0, 2, 0, 5, 5, 5, 5, 2, 3],
  [6, 0, 6, 2, 5, 5, 0, 3, 5, 5, 5, 5, 2],
  [5, 4, 0, 0, 2, 6, 0, 2, 3, 3, 5, 5, 5],
  [2, 5, 1, 5, 4, 0, 0, 3, 3, 5, 5, 5, 5],
  [6, 6, 2, 0, 4, 5, 4, 3, 4, 1, 0, 6, 4],
  [4, 4, 2, 5, 0, 3, 4, 3, 3, 4, 1, 0, 6],
  [5, 2, 0, 3, 2, 0, 4, 1, 0, 6, 4, 3, 3],
  [3, 0, 6, 3, 0, 6, 4, 4, 3, 3, 4, 1, 0],
  [3, 2, 0, 0, 6, 0, 4, 4, 1, 0, 6, 4, 3],
  [0, 0, 4, 6, 0, 4, 4, 6, 4, 3, 3, 4, 1],
  [6, 4, 2, 5, 6, 6, 1, 1, 4, 3, 4, 6, 3],
]
determinant_modulo_7 = 1
\end{verbatim}
\endgroup

\begin{thebibliography}{99}
\bibitem{sevenFukuda1994}
Fukuda, T. (1994).
Remarks on $\mathbb Z_p$-extensions of number fields.
\emph{Proceedings of the Japan Academy, Series A, Mathematical Sciences},
\emph{70}(8), 264--266.
\url{https://doi.org/10.3792/pjaa.70.264}.

\bibitem{sevenFriedman1989}
Friedman, E. (1989).
Analytic formulas for the regulator of a number field.
\emph{Inventiones Mathematicae}, \emph{98}(3), 599--622.
\url{https://doi.org/10.1007/BF01393839}.

\bibitem{sevenSutherland2015}
Sutherland, A. V. (2015, November 12).
\emph{The analytic class number formula} (18.785 Lecture Notes 18).
Massachusetts Institute of Technology.
\url{https://math.mit.edu/classes/18.785/2015fa/LectureNotes18.pdf}.

\bibitem{sevenSimard2015}
Simard, N. (2015, November 2).
\emph{Notes on theta series}.
\url{https://www.math.mcgill.ca/goren/GS/Thetas.pdf}.

\bibitem{sevenDarmonLauderRotger2015}
Darmon, H., Lauder, A., \& Rotger, V. (2015).
Stark points and $p$-adic iterated integrals attached to modular forms of
weight one.
\emph{Forum of Mathematics, Pi}, \emph{3}, e8.
\url{https://doi.org/10.1017/fmp.2015.7}.

\bibitem{sevenJohansson2017}
Johansson, F. (2017).
Arb: Efficient arbitrary-precision midpoint-radius interval arithmetic.
\emph{IEEE Transactions on Computers}, \emph{66}(8), 1281--1292.
\url{https://doi.org/10.1109/TC.2017.2690633}.

\bibitem{sevenSharifi}
Sharifi, R. (n.d.).
\emph{Iwasawa theory} [Author's lecture notes].
\url{https://www.math.ucla.edu/~sharifi/notes/iwasawa.pdf}.
Chapters 1 and 3, especially Corollaries 1.3.9 and 1.3.12,
Remark 1.3.21, and Proposition 3.2.20.

\bibitem{Rubin} Rubin, K. (1991). The ``main conjectures'' of Iwasawa theory for imaginary quadratic fields. \emph{Inventiones Mathematicae, 103}, 25--68. \url{https://doi.org/10.1007/BF01239508}.
\bibitem{Viguie} Vigui\'e, S. (2011). \emph{On the classical main conjecture for imaginary quadratic fields}. arXiv:1103.1125v4, Theorem 1.1 and Sections 1--3. \url{https://arxiv.org/abs/1103.1125}.
\end{thebibliography}
\end{document}
